\documentclass[11pt,a4paper]{article}

% --- Paket LaTeX Standar ---
\usepackage[utf8]{inputenc}
\usepackage[bahasa]{babel}
\usepackage[margin=1in]{geometry}
\usepackage{amsmath,amsfonts,amssymb}
\usepackage{booktabs}
\usepackage{xcolor}
\usepackage{graphicx}
\usepackage{listings}
\usepackage{enumitem}
\usepackage{fancyhdr}
\usepackage{tcolorbox}
\usepackage{hyperref}
\usepackage{mathptmx}      % Menggunakan font Times New Roman
\usepackage{parskip}       % Spasi antar paragraf modern & tanpa indentasi paragraf awal
\renewcommand{\arraystretch}{1.3} % Membuat tinggi baris tabel lebih renggang & profesional

% --- Konfigurasi Hyperref ---
\hypersetup{
    colorlinks=true,
    linkcolor=blue,
    filecolor=magenta,      
    urlcolor=cyan,
    pdftitle={Rencana Anggaran Biaya (RAB) - Sistem ERP PT Semadam},
    pdfpagemode=FullScreen,
}

% --- Konfigurasi Warna & Kotak Pesan ---
\definecolor{primarycolor}{HTML}{0f172a} % Slate 900
\definecolor{secondarycolor}{HTML}{ea580c} % Orange 600
\definecolor{alertbg}{HTML}{f8fafc}
\definecolor{alertborder}{HTML}{cbd5e1}
\definecolor{alerttext}{HTML}{334155}

\newtcolorbox{notealert}[1][Catatan Finansial]{
    colback=alertbg,
    colframe=alertborder,
    coltext=alerttext,
    fonttitle=\bfseries,
    title=#1,
    arc=4px,
    boxrule=1pt,
    left=10pt, right=10pt, top=8pt, bottom=8pt
}

% --- Konfigurasi Header & Footer ---
\pagestyle{fancy}
\fancyhf{}
\fancyhead[L]{\small\color{gray}Sistem ERP PT Semadam - Rencana Anggaran Biaya (RAB)}
\fancyhead[R]{\small\color{gray}Alokasi Nilai Fitur}
\fancyfoot[L]{\small\color{gray}\copyright~2026 PT Semadam}
\fancyfoot[R]{\small\thepage}
\renewcommand{\headrulewidth}{0.4pt}
\renewcommand{\footrulewidth}{0.4pt}

% --- INFORMASI DOKUMEN ---
\title{
    \vspace{-1.5cm}
    \color{primarycolor}\textbf{\Huge RENCANA ANGGARAN BIAYA (RAB)}\\
    \vspace{0.3cm}
    \color{secondarycolor}\Large\textbf{Distribusi Anggaran Fitur \& Milestones Finansial ERP}\\
    \vspace{0.2cm}
    \large\textbf{Sistem Informasi ERP Terintegrasi PT Semadam}
}
\author{\textbf{PT Semadam - Tim Pengembang Sistem}}
\date{29 Mei 2026}

\begin{document}

\maketitle

\thispagestyle{empty}
\begin{center}
    \vspace{1.5cm}
    \begin{tabular}{ll}
        \toprule
        \textbf{Detail Dokumen} & \textbf{Informasi} \\
        \midrule
        Nama Proyek & Sistem ERP Terintegrasi PT Semadam \\
        Dokumen / Lampiran & Rencana Anggaran Biaya (RAB) \& Alokasi Harga Fitur \\
        Versi Dokumen & 1.0.0 (Final) \\
        Status Dokumen & Lampiran Kontrak Resmi (Disetujui) \\
        Bahasa & Bahasa Indonesia \\
        \bottomrule
    \end{tabular}
\end{center}

\newpage
\tableofcontents
\newpage

% ==========================================
\section{PENDAHULUAN}
% ==========================================

\subsection{Latar Belakang}
Pengembangan Sistem Informasi ERP Terintegrasi PT Semadam mencakup lima modul operasional dan finansial terpadu yang berjalan secara paralel. Untuk memastikan transparansi nilai pekerjaan dan menyusun milestones pembayaran berbasis fungsionalitas (\textit{feature-based payment milestones}), diperlukan perincian anggaran biaya untuk setiap fitur di masing-masing modul.

\subsection{Metode Alokasi Anggaran}
Dalam menetapkan nilai nominal per fitur, tim pengembang tidak menerapkan pembagian rata (\textit{flat division}), melainkan menggunakan \textbf{Metode Pembobotan Kompleksitas Fitur (Complexity-Weighted Allocation)}. Metode ini dinilai jauh lebih profesional dan akurat karena beberapa alasan:
\begin{enumerate}
    \item \textbf{Representasi Usaha (Effort Representation)}: Fitur yang membutuhkan arsitektur database rumit, Server Actions mutakhir, dan algoritma komputasi berat (seperti \textit{GL Engine} dan \textit{FIFO Valuation}) dinilai dengan harga lebih tinggi dibanding CRUD master data dasar.
    \item \textbf{Keadilan Finansial (Financial Fairness)}: Menyelaraskan harga fitur dengan keahlian teknis (\textit{expertise}) yang dibutuhkan untuk membangunnya.
    \item \textbf{Kemudahan Audit}: Memudahkan manajemen PT Semadam dalam melakukan audit progress serah terima pekerjaan per fitur secara objektif.
\end{enumerate}

\newpage

% ==========================================
\section{RINGKASAN ANGGARAN MODUL ERP}
% ==========================================

Total nilai kontrak pengembangan Sistem ERP Terintegrasi PT Semadam adalah \textbf{Rp 92.000.000,- (Sembilan Puluh Dua Juta Rupiah)} dengan rincian per modul sebagai berikut:

\begin{table}[h!]
\centering
\small
\begin{tabular}{llrlc}
    \toprule
    \textbf{No.} & \textbf{Modul ERP} & \textbf{Harga Modul (IDR)} & \textbf{Persentase (\%)} & \textbf{Status Pengembangan} \\
    \midrule
    1 & \textbf{Accounting (SIA)} & Rp 20.000.000,- & 21,74\% & Financial Core Hub \\
    2 & \textbf{Kas/Bank (Treasury)} & Rp 15.000.000,- & 16,30\% & Cash Flow Engine \\
    3 & \textbf{Inventory (Gudang)} & Rp 27.000.000,- & 29,35\% & Logistics \& Stocks \\
    4 & \textbf{Gaji (Payroll)} & Rp 17.000.000,- & 18,48\% & Plantation HRD Engine \\
    5 & \textbf{Assets (Aktiva Tetap)} & Rp 15.000.000,- & 16,13\% & Depreciation Engine \\
    \midrule
    & \textbf{TOTAL} & \textbf{Rp 92.000.000,-} & \textbf{100,00\%} & \textbf{Terintegrasi Penuh} \\
    \bottomrule
\end{tabular}
\caption{Ringkasan Anggaran Pengembangan Modul ERP PT Semadam}
\end{table}

\begin{notealert}[Integrasi Finansial]
Meskipun harga dibagi per modul, kelima modul di atas dirancang untuk beroperasi secara terintegrasi penuh melalui backend PostgreSQL Supabase secara \textit{real-time} tanpa penundaan sinkronisasi.
\end{notealert}

\newpage

% ==========================================
\section{DETAIL ALOKASI ANGGARAN FITUR PER MODUL}
% ==========================================

\subsection{Modul 1: Accounting (SIA) - Total Rp 20.000.000,-}
Modul ini bertindak sebagai jantung keuangan konsolidasi seluruh unit kebun PT Semadam.

\begin{table}[h!]
\centering
\scriptsize
\begin{tabular}{llcrl}
    \toprule
    \textbf{Nama Fitur / Komponen} & \textbf{Kompleksitas} & \textbf{Bobot} & \textbf{Nilai Fitur (IDR)} & \textbf{Deskripsi Fungsional} \\
    \midrule
    \textbf{Core GL \& Trial Balance} & Tinggi & 30\% & Rp 6.000.000,- & Server Action \texttt{calculateTrialBalance} (Saldo Awal \& Mutasi) \\
    \textbf{LNET \& Laporan Manajemen} & Tinggi & 27,5\% & Rp 5.500.000,- & Komparasi Aktual vs Anggaran, cache tabel \texttt{laporan\_manajemen\_netral} \\
    \textbf{Memorial Journal Form} & Sedang & 17,5\% & Rp 3.500.000,- & Form input berimbang, autokomplit COA, hotkey \texttt{Ctrl+S} \\
    \textbf{Master Data CRUD} & Rendah & 15\% & Rp 3.000.000,- & CRUD \texttt{master\_unit}, \texttt{master\_rekening}, Excel Import COA \\
    \textbf{Floating Calculator \& Console} & Sedang & 10\% & Rp 2.000.000,- & Widget kalkulator melayang \& grid admin perbaikan darurat \\
    \midrule
    \textbf{TOTAL} & & \textbf{100\%} & \textbf{Rp 20.000.000,-} & \\
    \bottomrule
\end{tabular}
\end{table}

\subsection{Modul 2: Kas/Bank (Treasury) - Total Rp 15.000.000,-}
Mengelola aliran masuk dan keluar dana cair kebun serta rekonsiliasi perbankan.

\begin{table}[h!]
\centering
\scriptsize
\begin{tabular}{llcrl}
    \toprule
    \textbf{Nama Fitur / Komponen} & \textbf{Kompleksitas} & \textbf{Bobot} & \textbf{Nilai Fitur (IDR)} & \textbf{Deskripsi Fungsional} \\
    \midrule
    \textbf{Cash Flow Ledger Engine} & Tinggi & 33,3\% & Rp 5.000.000,- & Otomasi posting jurnal transaksi kas keluar/masuk ke Ledger \\
    \textbf{Bank Reconciliation} & Tinggi & 30\% & Rp 4.500.000,- & Komparasi mutasi kas internal vs rekening koran via CSV bank \\
    \textbf{Voucher Entry Forms} & Sedang & 23,3\% & Rp 3.500.000,- & UI pembuatan slip bukti kas masuk/keluar multi-baris pembebanan \\
    \textbf{Master Kas \& Rekening Bank} & Rendah & 13,4\% & Rp 2.000.000,- & Setup kasir kebun, limit harian, otorisasi tanda tangan slip \\
    \midrule
    \textbf{TOTAL} & & \textbf{100\%} & \textbf{Rp 15.000.000,-} & \\
    \bottomrule
\end{tabular}
\end{table}

\subsection{Modul 3: Inventory (Logistik \& Gudang) - Total Rp 27.000.000,-}
Modul terbesar menangani suku cadang pabrik sawit, pupuk, BBM, dan logistik kebun.

\begin{table}[h!]
\centering
\scriptsize
\begin{tabular}{llcrl}
    \toprule
    \textbf{Nama Fitur / Komponen} & \textbf{Kompleksitas} & \textbf{Bobot} & \textbf{Nilai Fitur (IDR)} & \textbf{Deskripsi Fungsional} \\
    \midrule
    \textbf{Stock Valuation Engine} & Tinggi & 31,5\% & Rp 8.500.000,- & Penilaian stok mutasi barang dengan metode \textit{Weighted Average} \\
    \textbf{Append Gudang \& Opname} & Tinggi & 24\% & Rp 6.500.000,- & Closing stok bulanan, stock opname fisik, pencatatan deviasi \\
    \textbf{Buku Penerimaan \& Pengeluaran} & Sedang & 22,2\% & Rp 6.000.000,- & Transaksi LPB/Nota dan pengeluaran bon barang per afdeling kebun \\
    \textbf{Master Barang \& Alert Order} & Sedang & 14,8\% & Rp 4.000.000,- & Setup batas minimal stok barang kebun \& alert restock otomatis \\
    \textbf{Supplier \& Purchase Order} & Rendah & 7,5\% & Rp 2.000.000,- & Setup rekanan suplier, cetak dokumen PO resmi, pelacakan transit \\
    \midrule
    \textbf{TOTAL} & & \textbf{100\%} & \textbf{Rp 27.000.000,-} & \\
    \bottomrule
\end{tabular}
\end{table}

\newpage

\subsection{Modul 4: Gaji (Payroll) - Total Rp 17.000.000,-}
Sistem manajemen ketenagakerjaan khusus perkebunan sawit (staf bulanan vs buruh panen).

\begin{table}[h!]
\centering
\scriptsize
\begin{tabular}{llcrl}
    \toprule
    \textbf{Nama Fitur / Komponen} & \textbf{Kompleksitas} & \textbf{Bobot} & \textbf{Nilai Fitur (IDR)} & \textbf{Deskripsi Fungsional} \\
    \midrule
    \textbf{Palm Oil Labor Payroll Engine} & Tinggi & 38,2\% & Rp 6.500.000,- & Hitung upah harian, premi panen (TBS/timbangan), lembur buruh \\
    \textbf{BPJS \& PPh 21 Calculator} & Sedang & 23,5\% & Rp 4.000.000,- & Potongan PPh 21 bulanan, iuran iuran BPJS TK \& Kesehatan otomatis \\
    \textbf{Master Karyawan \& Absensi} & Sedang & 20,6\% & Rp 3.500.000,- & DB SKU/BHL, denda keterlambatan, integrasi tracker finger print \\
    \textbf{Slip Gaji Generator \& Rekap} & Rendah & 17,7\% & Rp 3.000.000,- & Cetak slip per divisi, rekap biaya upah bulanan, ekspor bank transfer \\
    \midrule
    \textbf{TOTAL} & & \textbf{100\%} & \textbf{Rp 17.000.000,-} & \\
    \bottomrule
\end{tabular}
\end{table}

\subsection{Modul 5: Assets (Aktiva Tetap) - Total Rp 15.000.000,-}
Pencatatan aset perkebunan, armada angkutan TBS, serta depresiasi aset berkala.

\begin{table}[h!]
\centering
\scriptsize
\begin{tabular}{llcrl}
    \toprule
    \textbf{Nama Fitur / Komponen} & \textbf{Kompleksitas} & \textbf{Bobot} & \textbf{Nilai Fitur (IDR)} & \textbf{Deskripsi Fungsional} \\
    \midrule
    \textbf{Automatic Depreciation Engine} & Tinggi & 36,6\% & Rp 5.500.000,- & Posting penyusutan bulanan otomatis (garis lurus/saldo menurun) \\
    \textbf{Asset Registry \& Lifecycle} & Sedang & 23,3\% & Rp 3.500.000,- & Track mutasi aset antar divisi kebun, perbaikan, disposal aset \\
    \textbf{CAPEX \& Asset Acquisition} & Sedang & 23,3\% & Rp 3.500.000,- & Alur persetujuan belanja modal pembelian alat berat/lahan baru \\
    \textbf{Master Kategori Asset} & Rendah & 16,8\% & Rp 2.500.000,- & Pengelompokan masa manfaat aset dan cetak barcode inventaris \\
    \midrule
    \textbf{TOTAL} & & \textbf{100\%} & \textbf{Rp 15.000.000,-} & \\
    \bottomrule
\end{tabular}
\end{table}

\newpage

% ==========================================
\section{METODE PEMBAYARAN \& TERMIN (PAYMENT MILESTONES)}
% ==========================================

Untuk menyelaraskan antara progress pengembangan fisik dengan komitmen keuangan, pembayaran dibagi menjadi \textbf{5 Termin} berdasarkan penyelesaian serah terima modul operasional:

\begin{enumerate}
    \item \textbf{Termin 1 (Uang Muka / DP) - 20\%} \\
    Pembayaran awal setelah penandatanganan \textit{Scope of Work} (SOW) dan persetujuan formal \textit{Product Requirement Document} (PRD). Nilai: \textbf{Rp 18.400.000,-}.
    \item \textbf{Termin 2 (Modul Accounting \& Kas/Bank) - 30\%} \\
    Setelah serah terima fungsionalitas penuh modul Keuangan \& Treasury (SIA \& Kas/Bank) serta lulus uji akurasi jurnal. Nilai: \textbf{Rp 27.600.000,-}.
    \item \textbf{Termin 3 (Modul Inventory / Gudang) - 25\%} \\
    Setelah modul persediaan barang, stok gudang, dan logistik kebun dinyatakan selesai uji fungsi dan *Go-Live*. Nilai: \textbf{Rp 23.000.000,-}.
    \item \textbf{Termin 4 (Modul Payroll \& Assets) - 20\%} \\
    Setelah modul rekap HRD/Gaji buruh panen dan modul depresiasi otomatis aktiva tetap berhasil diserahterimakan. Nilai: \textbf{Rp 18.400.000,-}.
    \item \textbf{Termin 5 (Pemeliharaan / Retensi) - 5\%} \\
    Dibayarkan setelah berakhirnya masa garansi/pemeliharaan gratis sistem selama 3 bulan operasional. Nilai: \textbf{Rp 4.600.000,-}.
\end{enumerate}

\vspace{2cm}

\begin{center}
    \rule{8cm}{0.5pt}\\
    \small\color{gray}DOKUMEN RENCANA ANGGARAN BIAYA (RAB) INI MERUPAKAN LAMPIRAN RESMI HUKUM YANG MERUPAKAN BAGIAN INTEGRAL DARI SCOPE OF WORK (SOW) PENGEMBANGAN ERP PT SEMADAM.
\end{center}

\end{document}
